\section{Low-height clustering in square blocks}
\label{sec:clustering}
%======================================================================

Low height is the part of the problem that geometry alone can control.  Here bounded factor gap becomes bounded occupancy, which yields a genuine theorem before any high-sector cancellation is assumed.

The $n$th square block becomes the vertical strip
\begin{equation}
 \mathscr B_n
 :=\{w_m:n^2\le\Repart w_m<(n+1)^2\}.
 \label{eq:complex-block}
\end{equation}
For a canonical pair $(c,q)$, define the absolute factor gap
\begin{equation}
 d:=|q-c|.
 \label{eq:absolute-gap}
\end{equation}
Then
\begin{equation}
 |Y|
 =\frac{|q^2-c^2|}{2}
 =\frac{d(q+c)}{2}.
 \label{eq:Y-absolute-gap}
\end{equation}
For a squarefree source $m>1$, one has $d\ge1$: equality $d=0$ would force $m=c^2$.

\begin{lemma}[At most one product per positive absolute gap]
\label{lem:one-product-gap}
Fix $n\ge1$ and an integer $d\ge1$.  There is at most one unordered positive factor pair $\{c,q\}$ with $|q-c|=d$ and
\begin{equation}
 n^2\le cq<(n+1)^2.
 \label{eq:gap-block-condition}
\end{equation}
\end{lemma}

\begin{proof}
Write the smaller factor as $r$ and the larger as $r+d$.  The product is $f_d(r)=r(r+d)$.  If $f_d(r)\in B_n$, then
\[
 2r+d\ge2\sqrt{r(r+d)}\ge2n.
\]
Hence
\[
 f_d(r+1)-f_d(r)=2r+d+1\ge2n+1,
\]
which exceeds the integer diameter $2n$ of the block.  Thus two distinct values of $r$ cannot both occur.
\end{proof}

\begin{theorem}[Sharpened absolute low-height clustering]
\label{thm:absolute-clustering}
For every $H\ge0$,
\begin{equation}
 \boxed{
 \#\{w_m\in\mathscr B_n:|Y_m|\le H\}
 \le
 \left\lfloor\frac{H}{n}\right\rfloor.
 }
 \label{eq:absolute-clustering}
\end{equation}
Moreover,
\begin{equation}
 \boxed{
 \#\{w_m\in\mathscr B_n:|Y_m|\le n\}=0.
 }
 \label{eq:no-height-n-points}
\end{equation}
\end{theorem}

\begin{proof}
If $cq\in B_n$, then
\[
 c+q\ge2\sqrt{cq}\ge2n.
\]
Therefore
\[
 |Y|=\frac{d(c+q)}2\ge dn.
\]
The condition $|Y|\le H$ implies $d\le\lfloor H/n\rfloor$.  Since $d\ge1$, there are at most $\lfloor H/n\rfloor$ possible gaps, and \cref{lem:one-product-gap} gives at most one point for each gap.

If $|Y|\le n$, the inequalities force $d=1$ and $c+q=2n$.  But $|q-c|=1$ makes $c+q$ odd, whereas $2n$ is even.  Hence equality is impossible and the set is empty.
\end{proof}

\subsection{Parity-refined bounds}

The factor gap has a prescribed parity:
\begin{itemize}[leftmargin=1.7em]
\item if $m$ is odd, then $c$ and $q$ are odd, so $d$ is even;
\item if $m$ is even and squarefree, then the canonical factors have opposite parity, so $d$ is odd.
\end{itemize}
Let
\begin{equation}
 D:=\left\lfloor\frac{H}{n}\right\rfloor.
 \label{eq:D-height}
\end{equation}
The permitted positive gaps are therefore
\[
 2,4,\ldots,2\left\lfloor\frac D2\right\rfloor
\]
for odd sources and
\[
 1,3,\ldots,2\left\lceil\frac D2\right\rceil-1
\]
for even sources.  Consequently,
\begin{align}
 \#\{w_m\in\mathscr B_n:m\text{ odd},\ |Y_m|\le H\}
 &\le
 \left\lfloor\frac{H}{2n}\right\rfloor,
 \label{eq:odd-parity-clustering}\\
 \#\{w_m\in\mathscr B_n:m\text{ even},\ |Y_m|\le H\}
 &\le
 \left\lceil\frac{1}{2}
 \left\lfloor\frac{H}{n}\right\rfloor\right\rceil.
 \label{eq:even-parity-clustering}
\end{align}
At the endpoint $H=n$, both populations are empty by \cref{eq:no-height-n-points}.

At height $H=\Lambda n$, the scale-free bounds become
\begin{align}
 \#\{w_m\in\mathscr B_n:|Y_m|\le\Lambda n\}
 &\le\lfloor\Lambda\rfloor,
 \label{eq:constant-occupancy}\\
 \#\{m\text{ odd}:|Y_m|\le\Lambda n\}
 &\le\left\lfloor\frac{\Lambda}{2}\right\rfloor.
 \label{eq:constant-odd-occupancy}
\end{align}

\subsection{A sharper low-height space--time budget}

For negative-height points, the transport lifetime is zero.  For positive-height points in $\mathscr B_n$, one has $q\ge\sqrt m\ge n$, so \cref{eq:L-Y} gives
\begin{equation}
 L(m)\le\frac{2H}{n}
 \qquad(0<Y_m\le H).
 \label{eq:L-H-n}
\end{equation}
Therefore
\begin{equation}
 \boxed{
 \sum_{\substack{w_m\in\mathscr B_n\\|Y_m|\le H}}
 (1+L(m))
 \le
 \left\lfloor\frac{H}{n}\right\rfloor
 \left(1+\frac{2H}{n}\right).
 }
 \label{eq:space-time-budget}
\end{equation}
The parity-refined versions are
\begin{align}
 \sum_{\substack{w_m\in\mathscr B_n\\m\text{ odd},\ |Y_m|\le H}}
 (1+L(m))
 &\le
 \left\lfloor\frac{H}{2n}\right\rfloor
 \left(1+\frac{2H}{n}\right),
 \label{eq:odd-space-time-budget}\\
 \sum_{\substack{w_m\in\mathscr B_n\\m\text{ even},\ |Y_m|\le H}}
 (1+L(m))
 &\le
 \left\lceil\frac{1}{2}
 \left\lfloor\frac{H}{n}\right\rfloor\right\rceil
 \left(1+\frac{2H}{n}\right).
 \label{eq:even-space-time-budget}
\end{align}
At $H=\Lambda n$, all three right sides are bounded independently of $n$.

%======================================================================
